% Generated draft: requires mathematical and source review.
\documentclass[11pt]{article}
\usepackage[margin=25mm]{geometry}
\usepackage{fontspec}
\setmainfont{Latin Modern Roman}
\usepackage{amsmath,amssymb,mathtools,unicode-math}
\setmathfont{Latin Modern Math}
\usepackage{enumitem,needspace,xurl,hyperref}
\setlength{\parindent}{0pt}
\setlength{\parskip}{5pt}
\setlength{\emergencystretch}{3em}
\newcommand{\N}{\mathbb{N}}
\newcommand{\Z}{\mathbb{Z}}
\newcommand{\Q}{\mathbb{Q}}
\newcommand{\R}{\mathbb{R}}
\newcommand{\C}{\mathbb{C}}
\newcommand{\F}{\mathbb{F}}
\newcommand{\A}{\mathbb{A}}
\newcommand{\PP}{\mathbb{P}}
\newcommand{\E}{\mathbb{E}}
\newcommand{\eps}{\varepsilon}
\newcommand{\dd}{\,\mathrm{d}}
\begin{document}
\section*{UnsolvedMath ID: 30003730}
\textbf{Source problem code:} catalogue entry
\subsection{Definitions and mathematical statement}
A random conductance model on \(\Z^d\) is a stochastic process whose jump rates are determined by a random environment of conductances. A quenched invariance principle means convergence in distribution of suitably rescaled processes under the process law with the environment fixed, for almost every environment with respect to its probability law.

\textbf{Problem.} For random conductance models with stable-like jumps on \(\Z^d\), determine the optimal integrability or moment condition replacing the sufficient bounds on positive and negative moments of the conductances that ensure the quenched invariance principle.

More explicitly, the desired conclusion has the schematic form
\[
\operatorname{Law}_{\omega}\!\left(X^{(\varepsilon)}\right)
\Longrightarrow \operatorname{Law}(Y)
\qquad\text{as }\varepsilon\downarrow0,
\quad\text{for }\mathbb P\text{-almost every }\omega,
\]
where \(\mathbb P\) denotes the environment law, \(X^{(\varepsilon)}\) denotes the appropriately rescaled process in that environment, and \(Y\) denotes the intended limiting process. These symbols describe the meaning of the requested quenched conclusion; the supplied record does not specify their model-dependent details.

\textbf{Unresolved qualifications.} The record does not specify the admissible dimensions, stable index, jump-rate kernel, environment assumptions, treatment of zero conductances, time parametrization, rescaling, limiting process, or convergence topology. It also omits the existing sufficient moment inequalities and the precise meaning of optimality. Consequently, it identifies a research target but does not determine a fully quantified sharp-condition problem. Those specifications must be recovered from the cited sources before a precise necessary-and-sufficient criterion or sharp range of moment exponents can be formulated.
\subsection{Short English statement}
Determine the sharp integrability assumptions on conductances that yield a quenched invariance principle for random conductance models with stable-like jumps on \(\Z^d\). The supplied record leaves the model assumptions, existing moment bounds, and meaning of sharpness unspecified; this is a statement-only draft requiring mathematical review.
\subsection{Sources}
\begin{itemize}
\item Supplied local UnsolvedMath registry record: permanent ID \(30003730\), source problem code \texttt{catalogue entry}, title \emph{Optimal Moment Conditions for Random Conductance Models}, in the Oberwolfach Reports Open Problems collection. The recorded proposed year is 2018; no individual proposer is supplied.
\item Supplied source citation: \emph{Interplay of Analysis and Probability in Applied Mathematics} (2019), Oberwolfach Reports, pp.\ 311--381, DOI \texttt{10.4171/owr/2018/6}. Supplied source URL: \url{https://doi.org/10.4171/owr/2018/6}.
\item Supplied literature source, labelled ``original OWR report'': \url{https://doi.org/10.4171/OWR/2018/6}.
\item Supplied literature source, labelled ``Chen--Kumagai--Wang, stable-like jumps'': \url{https://doi.org/10.1214/20-AAP1616}.
\item Literature assessment dated exactly \texttt{2026-08-21}: the supplied curation labels the problem open and reports that the cited stable-like random-conductance invariance principle uses explicit positive and negative moment assumptions described as far from optimal. This is supplied dated curation, not an independently verified literature assessment. The record expressly requests independent verification of time-sensitive status and cited claims.
\item The registry supplies the statement-status label \texttt{corrected\_verified} and reports manual semantic, typographic, and LaTeX review. These are supplied editorial assertions; they do not establish mathematical completeness or correctness of this draft.
\end{itemize}
\end{document}
